\section*{Bab \ref{ch:b1:electronic-effects} --- Efek Elektronik dan Zat Antara Reaktif}

\begin{solution}{exo:b1:electronic-effects:1}
2-Metilbutana, \ce{CH3-CH(CH3)-CH2-CH3}: ketiga \ce{CH3} primer, \ce{CH2}
sekunder, CH tersier. 2-Bromo-2-metilpropana: ketiga \ce{CH3} primer,
karbon tengahnya tersier. Dari keempat bromida \ce{C4H9Br}
(1-bromobutana, 2-bromobutana, 1-bromo-2-metilpropana,
2-bromo-2-metilpropana), yang terakhir menghasilkan \omterm{def:b1:electronic-effects:intermediates}{karbokation} tersier,
yang paling stabil.
\end{solution}

\begin{solution}{exo:b1:electronic-effects:2}
$10^{4.76 - 2.87} = 10^{1.89} = 78$; $10^{4.76 - 0.52} = 10^{4.24} =
\num{1.7e4}$.
\end{solution}

\begin{solution}{exo:b1:electronic-effects:3}
Etanoat: C=O dan \ce{C-O^-} dipertukarkan antara kedua oksigen.
4-Nitrofenoksida: \omterm{def:b1:lewis-resonance-vsepr:lewis}{pasangan elektron bebas} oksigen membentuk C=O,
ikatan-ikatan rangkap cincin bergeser, karbon yang membawa \ce{NO2}
membentuk ikatan C=N, dan sepasang elektron N=O berpindah ke oksigen:
struktur dengan C=O di satu ujung cincin, C=N di ujung lainnya, dan
kedua oksigen nitro bermuatan negatif (sebuah struktur kuinoid). Kation
alil: \ce{CH2=CH-CH2+} $\leftrightarrow$ \ce{+CH2-CH=CH2}.
\end{solution}

\begin{solution}{exo:b1:electronic-effects:4}
Panah kedua: dari ikatan O--H air ke oksigennya; produknya \ce{H2O}
(dari \ce{OH-}, kini netral) dan \ce{OH-} (oksigen bekas air
mempertahankan pasangan itu, muatan $-1$). Untuk amonia: panah dari
\omterm{def:b1:lewis-resonance-vsepr:lewis}{pasangan elektron bebas} N ke salah satu H pada \ce{H3O+}, dan satu panah
dari ikatan O--H itu ke oksigen: N menjadi $+1$, O menjadi 0.
\end{solution}

\begin{solution}{exo:b1:electronic-effects:5}
Etanol (15.9) $<$ air (14.0) $<$ fenol (9.99, delokalisasi dalam
cincin) $<$ 4-nitrofenol (7.16, \ce{NO2} mengambil muatannya) $<$ asam
etanoat (4.76, dua oksigen yang ekuivalen) $<$ asam trikloroetanoat
(0.51, tiga klor $-I$).
\end{solution}

\begin{solution}{exo:b1:electronic-effects:6}
Anilina (4.60) $<$ piridina (5.23) $<$ amonia (9.25) $<$ metilamina
(10.66). Pada anilina \omterm{def:b1:lewis-resonance-vsepr:lewis}{pasangan elektron bebasnya} dipakai bersama dengan
cincin (struktur-struktur dengan C=N$^+$ dan muatan negatif pada
\textit{orto} dan \textit{para}); memprotonasi N akan mengorbankan
delokalisasi itu.
\end{solution}

\begin{solution}{exo:b1:electronic-effects:7}
Pada 3-nitrofenoksida, struktur-struktur resonansi menempatkan muatan pada
karbon-karbon \textit{orto} dan \textit{para} terhadap oksigen, tidak
pernah pada karbon yang membawa \ce{NO2}: tidak ada bantuan mesomeri,
sehingga lebih lemah daripada isomer 4. Gugus nitro tetap menarik
elektron melalui ikatan-ikatan $\sigma$ ($-I$), sehingga lebih kuat
daripada fenol.
\end{solution}

\begin{solution}{exo:b1:electronic-effects:8}
\ce{CH3+} $<$ \ce{CH3CH2+} $<$ \ce{CH2=CH-CH2+} $<$ \ce{C6H5CH2+}
$\approx$ \ce{(CH3)3C+} $<$ \ce{CH3OCH2+}: gugus alkil memberikan
elektron; alil dan benzil mendelokalisasi muatannya (dua dan empat
struktur); \omterm{def:b1:lewis-resonance-vsepr:lewis}{pasangan elektron bebas} oksigen menghasilkan struktur dengan
oktet penuh di mana-mana. Urutan di bagian tengah daftar itu bergantung
pada mediumnya.
\end{solution}

\begin{solution}{exo:b1:electronic-effects:9}
\ce{OH-} (\omterm{def:b1:electronic-effects:nucleophile}{nukleofil}) menyerang karbon \ce{CH3Br} (\omterm{def:b1:electronic-effects:nucleophile}{elektrofil}), dengan
pasangan C--Br pergi bersama Br. \ce{H2O} memberikan sepasang elektron
kepada \ce{H+} (di sini nama yang dipakai adalah asam dan basa).
\ce{CN-} menyerang karbon karbonil etanal, dengan pasangan $\pi$ C=O
berpindah ke oksigen. \ce{NH3} memberikan \omterm{def:b1:lewis-resonance-vsepr:lewis}{pasangan elektron bebasnya}
kepada orbital kosong boron dalam \ce{BF3}.
\end{solution}

\begin{solution}{exo:b1:electronic-effects:10}
Tidak: $K_a = 10^{-0.51} = 0.31$ jauh lebih besar daripada $C$; asam itu
hampir terdisosiasi sempurna. $x^2/(0.010 - x) = 0.31$ menghasilkan
$x = \qty{9.7e-3}{mol/L}$, $\mathrm{pH} = 2.01$, \omterm{def:b1:predominant-reaction:dissociation}{derajat disosiasi}
\qty{97}{\percent}.
\end{solution}

\begin{solution}{exo:b1:electronic-effects:11}
Gugus-gugus alkil mendorong elektron ($+I$) ke oksigen alkoksida dan
membuat muatannya kurang stabil, makin banyak gugusnya makin besar
pengaruhnya. Alkoksida yang lebih besar juga kurang baik tersolvasi oleh
air, yang menstabilkan oksigen kecil bermuatan lebih baik daripada
oksigen yang terbenam di antara gugus-gugus alkil. Keduanya menaikkan
$\mathrm{p}K_a$.
\end{solution}

\begin{solution}{exo:b1:electronic-effects:12}
Dalam amida \omterm{def:b1:lewis-resonance-vsepr:lewis}{pasangan elektron bebas} nitrogen dipakai bersama dengan C=O
(struktur \ce{N+=C-O^-}); mengambilnya untuk sebuah proton atau \omterm{def:b1:electronic-effects:nucleophile}{elektrofil}
akan merusak stabilisasi itu. Dalam \omterm{def:b1:predominant-reaction:strong-acid}{asam kuat} amida terprotonasi pada
oksigennya, yang membawa muatan parsial negatif.
\end{solution}

\begin{solution}{pb:b1:electronic-effects:1}
\textbf{1.} $10^{1.89} = 78$; $10^{1.52} = 33$; $10^{0.84} = 6.9$.
\textbf{2.} Tidak: setiap klor tambahan membantu lebih sedikit. Muatan
karboksilat sudah tersebar dengan baik, dan pengukurannya menjadi kurang
teliti ketika asam itu mendekati disosiasi sempurna.
\textbf{3.} Setiap Cl menarik rapatan elektron melalui ikatan C--C dan
C--O serta menstabilkan karboksilat ($-I$).
\textbf{4.} \omterm{def:b1:electronic-effects:inductive}{Efek induktif} memudar dalam dua atau tiga ikatan.
\textbf{5.} Kedua asam itu tampak sama kuat (0.52 dan 0.51). Di dalam air,
keduanya hampir terdisosiasi sempurna pada konsentrasi yang biasa
(\cref{ch:b1:predominant-reaction}, perataan): $\mathrm{p}K_a$ di dekat 0
adalah batas yang dapat dibedakan oleh air, dan kedua nilai itu membawa
ketidakpastian beberapa persepuluh.
\textbf{6.} Asam etanoat: bentuk asam (\qty{98}{\percent}); kloroetanoat:
keduanya, dengan anion sedikit predominan ($10^{3 - 2.87} = 1.3$); ketiga
asam lainnya: anion.
\textbf{7.} \ce{CH3-C(=O)-O^-} $\leftrightarrow$ \ce{CH3-C(-O^-)=O}.
\textbf{8.} Dua ikatan C--O yang sama, di antara ikatan rangkap dua dan
ikatan tunggal.
\textbf{9.} Muatannya berada pada satu-satunya oksigen, tanpa ikatan
$\pi$ untuk membaginya.
\textbf{10.} $10^{15.9 - 4.76} = 10^{11.1}$, sekitar $10^{11}$.
\textbf{11.} Muatan fenoksida tersebar ke tiga karbon cincin, yang kurang
elektronegatif daripada oksigen: sebuah keuntungan, yang lebih kecil
daripada pada karboksilat: $\mathrm{p}K_a$ 9.99 di antara 4.76 dan 15.9.
\textbf{12.} $K = 10^{6.37 - \mathrm{p}K_a}$: asam etanoat $10^{1.61} = 41$,
disukai (karbon dioksida dilepaskan dengan hidrogenkarbonat berlebih);
fenol $10^{-3.62}$ dan etanol $10^{-9.5}$: tidak bereaksi.
\textbf{13.} 4-Nitrofenol (7.16) $>$ 2-nitrofenol (7.23) $>$ 3-nitrofenol
(8.36) $>$ fenol (9.99).
\textbf{14.} Struktur kuinoid pada \cref{exo:b1:electronic-effects:3},
dengan muatan pada salah satu oksigen gugus nitro.
\textbf{15.} Muatan itu hanya mencapai karbon-karbon \textit{orto} dan
\textit{para} terhadap oksigen; karbon nitro berada pada posisi
\textit{meta}.
\textbf{16.} Efek $-I$ gugus \ce{NO2}.
\textbf{17.} Pada 2-nitrofenol, OH membentuk \omterm{def:b1:intermolecular-forces-solvents:hbond}{ikatan hidrogen} dengan salah
satu oksigen gugus nitro di sebelahnya, yang menstabilkan bentuk asam
dan membuat protonnya sedikit lebih sukar dilepaskan.
\textbf{18.} Fraksi anion $1/(1 + 10^{\mathrm{p}K_a - 7.5})$: 4-nitro
\qty{69}{\percent}, 2-nitro \qty{65}{\percent}, 3-nitro
\qty{12}{\percent}. Bentuk asam dan anion 4-nitrofenol menyerap secara
berbeda, anionnya di daerah tampak: warnanya muncul di sekitar
$\mathrm{p}K_a$-nya.
\textbf{19.} $x^2/(0.010 - x) = 10^{-0.52} = 0.30$: $x = \qty{9.7e-3}{mol/L}$,
$\mathrm{pH} = 2.01$: hampir merupakan \omterm{def:b1:predominant-reaction:strong-acid}{asam kuat}.
\textbf{20.} $\frac12(4.76 + 2.00) = 3.38$.
\textbf{21.} \ce{CF3COOH + CH3COO- -> CF3COO- + CH3COOH}, $K = 10^{4.76 -
0.52} = 10^{4.24}$: \omterm{def:b1:extent-q-and-k:quantitative}{kuantitatif}.
\textbf{22.} $K_a(\ce{CF3COOH})/K_a(\ce{CH3COOH}) =$ \textbf{$10^{4.24} =
\num{1.7e4}$}, sekitar sepuluh ribu.
\end{solution}
